\documentclass{article}

\usepackage[margin=1in]{geometry}
\usepackage[T1]{fontenc}
\usepackage{lmodern}
\usepackage{amsmath,amssymb,amsthm,mathtools}
\usepackage{array,booktabs,longtable}
\usepackage[hidelinks]{hyperref}
\usepackage[nameinlink,noabbrev]{cleveref}
\usepackage{enumitem}

% LeanBlueprint-style commands
\providecommand{\lean}[1]{\texttt{#1}}
\providecommand{\leanok}{}
\providecommand{\mathlibok}{}
\providecommand{\uses}[1]{\par\smallskip\noindent\emph{Uses:} #1.\par\smallskip}
\providecommand{\proves}[1]{\par\smallskip\noindent\emph{Proves:} #1.\par\smallskip}
\providecommand{\notready}{\par\smallskip\noindent\emph{Status: not yet formalized.}\par\smallskip}

\newtheorem{definition}{Definition}[section]
\newtheorem{lemma}[definition]{Lemma}
\newtheorem{proposition}[definition]{Proposition}
\newtheorem{theorem}[definition]{Theorem}
\newtheorem{corollary}[definition]{Corollary}
\theoremstyle{remark}
\newtheorem{remark}[definition]{Remark}

\newcolumntype{L}[1]{>{\raggedright\arraybackslash}p{#1}}

\newcommand{\C}{\mathbb{C}}
\newcommand{\R}{\mathbb{R}}
\newcommand{\E}{\mathbb{E}}
\newcommand{\Prb}{\mathbb{P}}
\newcommand{\id}{\mathrm{Id}}
\newcommand{\Id}{\mathrm{Id}}
\newcommand{\tr}{\operatorname{tr}}
\newcommand{\Mat}{\operatorname{Mat}}
\newcommand{\code}{\mathcal{C}}
\newcommand{\hilb}{\mathcal{H}}
\newcommand{\algebra}{\mathcal{A}}
\newcommand{\eps}{\varepsilon}

\title{Strict-Proof-Following Blueprint for arXiv:2111.08131\\
\large\textit{Quantum Soundness of Testing Tensor Codes}}
\author{MIPStarRE Project}
\date{March 2026}

\begin{document}
\maketitle
\tableofcontents

%% ====================================================================
\section{Scope and design principles}
\label{sec:scope}

This document is a \emph{strict-proof-following} blueprint for formalizing
arXiv:2111.08131 (Natarajan--Vidick, ``Quantum soundness of testing tensor
codes''). It is \textbf{not} a pilot or a surrogate plan; it is intended to
mirror the proof in the paper as faithfully as possible.

\begin{description}[style=nextline,leftmargin=2.5cm]
	\item[Primary target]
	      \texttt{thm:main} (Theorem~4.1): quantum soundness of the tensor code test
	      for synchronous/tracial strategies on a von Neumann algebra with a normal
	      tracial state.
	\item[Secondary target]
	      \texttt{thm:main-bipartite} (Theorem~4.8): extension to general two-prover
	      (non-synchronous) strategies, using~\cite{vidick2021almost}.
	\item[Ambient setting]
	      Throughout, $\algebra$ is a von Neumann algebra on a separable Hilbert space
	      $\hilb$, and $\tau\colon\algebra\to\C$ is a normal tracial state.
	      All measurements take values in $\algebra$. This is the paper's ambient
	      setting; it is \textbf{not} replaced by a finite-dimensional matrix pilot.
	\item[Proof order]
	      The blueprint follows the paper's proof order:
	      Section~2 (preliminaries) $\to$ Section~3 (test definition) $\to$
	      Appendix~A (expander) $\to$ Section~5 (self-improvement) $\to$
	      Section~6 (pasting) $\to$ Section~4 (main induction, main theorems).
\end{description}

%% ====================================================================
\section{Section~2: Preliminaries}
\label{sec:prelims}

\subsection{Von Neumann algebras and tracial states}

The paper works in the following setting.

\begin{definition}[Normal tracial state]\label{def:tracial-state}
	Let $\algebra\subseteq\mathcal{B}(\hilb)$ be a von Neumann algebra.
	A \emph{normal tracial state} is a positive linear functional
	$\tau\colon\algebra\to\C$ that is unital ($\tau(\id)=1$),
	normal (represented by a positive trace-class operator),
	and tracial ($\tau(AB)=\tau(BA)$ for all $A,B\in\algebra$).
\end{definition}

\begin{definition}[$\tau$-norm]\label{def:tau-norm}
	\uses{def:tracial-state}
	For $A\in\algebra$,
	$\|A\|_\tau = \sqrt{\tau(A^*A)}$.
\end{definition}

\begin{proposition}[Triangle and H\"older inequalities --- \texttt{prop:holder}]
	\label{prop:holder}
	\uses{def:tracial-state, def:tau-norm}
	Let $\tau$ be a tracial state on a von Neumann algebra $\algebra$
	and let $A,B\in\algebra$. Then:
	\begin{enumerate}[label=(\roman*)]
		\item\label{item:holder-0} $|\tau(A)|\le\tau(|A|)$.
		\item\label{item:holder-1} (H\"older~1) $\tau(|AB|)\le\|A\|_\tau\cdot\|B\|_\tau$.
		\item\label{item:holder-2} (H\"older~2) $\tau(|AB|)\le\|A\|\cdot\tau(|B|)$.
		\item\label{item:triangle} (Triangle for $1$-norm) $\tau(|A+B|)\le\tau(|A|)+\tau(|B|)$.
	\end{enumerate}
\end{proposition}

\subsection{Measurements, consistency, and closeness}

\begin{definition}[Submeasurement, measurement, projective measurement]\label{def:submeasurement}
	\uses{def:tracial-state}
	Let $\mathcal{A}$ be a finite set.
	A \emph{submeasurement} on $\mathcal{A}$ is a family
	$\{M_a\}_{a\in\mathcal{A}}\subset\algebra$ with each $M_a\ge 0$ and
	$\sum_a M_a\le\id$.
	It is a \emph{measurement} if $\sum_a M_a=\id$,
	and \emph{projective} if each $M_a$ is a projection.
\end{definition}

\begin{definition}[Processed measurement]\label{def:processed}
	\uses{def:submeasurement}
	If $f\colon\mathcal{A}\to\mathcal{B}$ is a map, the processed submeasurement
	$M_{[f\mid b]}=\sum_{a:f(a)=b}M_a$.
\end{definition}

\begin{definition}[Consistency and closeness]\label{def:consistency}
	\uses{def:submeasurement, def:tau-norm}
	Let $\mu$ be a distribution on a finite question set $\mathcal{X}$
	and let $M^x,N^x$ be submeasurements indexed by $x\in\mathcal{X}$
	with a common finite answer set $\mathcal{A}$.
	\begin{itemize}
		\item $M^x_a\simeq_\delta N^x_a$ ($\delta$-consistent) iff
		      $\E_{x\sim\mu}\sum_{a\ne b}\tau(M^x_a N^x_b)\le\delta$.
		\item $M^x_a\approx_\delta N^x_a$ ($\delta$-close) iff
		      $(\E_{x\sim\mu}\sum_a\|M^x_a-N^x_a\|_\tau^2)^{1/2}\le\delta$.
	\end{itemize}
\end{definition}

\begin{proposition}[Cauchy--Schwarz for $\tau$ --- \texttt{cor:cauchy-schwarz}]
	\label{cor:cauchy-schwarz}
	\uses{def:tracial-state}
	$|\tau(A^*B)|^2\le\tau(A^*A) \tau(B^*B)$ for all $A,B\in\algebra$.
\end{proposition}

\begin{proposition}[Data processing --- \texttt{lem:data-processing}]
	\label{lem:data-processing}
	\uses{def:processed, def:consistency}
	If $M^x_a\simeq_\delta N^x_a$ and $f\colon\mathcal{A}\to\mathcal{B}$,
	then $M^x_{[f\mid b]}\simeq_\delta N^x_{[f\mid b]}$.
\end{proposition}

\begin{proposition}[Consistency to closeness --- \texttt{lem:consistency-consequences}]
	\label{lem:consistency-consequences}
	\uses{def:consistency}
	If $M^x$ and $N^x$ are measurements and $M^x_a\simeq_\delta N^x_a$,
	then $M^x_a\approx_{\sqrt{2\delta}} N^x_a$.
\end{proposition}

\begin{proposition}[Closeness to consistency --- \texttt{lem:closeness-to-consistency}]
	\label{lem:closeness-to-consistency}
	\uses{def:consistency}
	If $M^x$ is a projective submeasurement and $N^x$ a submeasurement with
	$M^x_a\approx_\delta N^x_a$, then $M^x_a\simeq_\delta N^x_a$.
\end{proposition}

\begin{proposition}[Closeness preserves inner products --- \texttt{lem:closeness-to-close-ips}]
	\label{lem:closeness-to-close-ips}
	\uses{def:consistency}
	\notready
	If $A^x_a\approx_\delta B^x_a$ and $\sum_a M^x_a(M^x_a)^*\le\id$ for
	every~$x$, then
	$\E_x\sum_a\tau(M^x_a A^x_a)\approx_\delta\E_x\sum_a\tau(M^x_a B^x_a)$.
\end{proposition}

\begin{proposition}[Multiplication preserves closeness --- \texttt{lem:add-a-proj}]
	\label{lem:add-a-proj}
	\uses{def:consistency}
	\notready
	If $M_a\approx_\delta N_a$ and $\{R_b\}$ satisfies $\sum_b R_b^*R_b\le\id$
	(or $\sum_b R_b R_b^*\le\id$), then left (resp.\ right) multiplication by the
	$R_b$ preserves the same closeness bound.
\end{proposition}

\begin{proposition}[Transfer consistency --- \texttt{lem:transfer-cons}]
	\label{lem:transfer-cons}
	\uses{def:consistency, lem:closeness-to-close-ips}
	\notready
	If $A^x_a\simeq_\delta C^x_a$ and $A^x_a\approx_\eps B^x_a$,
	then $B^x_a\simeq_{\delta+\eps} C^x_a$.
\end{proposition}

\begin{proposition}[Consistency controls a submeasurement --- \texttt{lem:cons-sub-meas}]
	\label{lem:cons-sub-meas}
	\uses{def:consistency}
	\notready
	If $A^x$ is a submeasurement, $B^x$ a measurement, and
	$A^x_a\simeq_\gamma B^x_a$, then
	$A^x_a\approx_{\sqrt\gamma} A^x_a B^x_a$
	and $A^x_a\approx_{2\sqrt\gamma} A^x B^x_a$,
	where $A^x=\sum_a A^x_a$.
\end{proposition}

\begin{proposition}[Switcheroo --- \texttt{lem:switcheroo}]
	\label{lem:switcheroo}
	\uses{lem:add-a-proj}
	\notready
	If $A^x_a A^y_b\approx_\eps A^y_b A^x_a$ on average and $P^s_{\vec a}$ is
	a length-$k$ product, then
	$P^s_{\vec a}A^x_b\approx_{k\eps}A^x_b P^s_{\vec a}$.
\end{proposition}

\subsection{Codes and tensor codes}

\begin{definition}[Linear code --- \texttt{def:code}]\label{def:code}
	An $[n,k,d]_\Sigma$ code is a $k$-dimensional linear subspace
	$\code\subseteq\Sigma^n$ with minimum distance~$d$.
\end{definition}

\begin{proposition}[Uniqueness from distance --- \texttt{prop:distance0}]
	\label{prop:distance0}
	\uses{def:code}
	With $t=n-d+1$, any $t$ distinct coordinates determine at most one codeword.
\end{proposition}

\begin{definition}[Interpolable code --- \texttt{def:interpolable}]\label{def:interpolable}
	\uses{def:code, prop:distance0}
	$\code$ is interpolable if for every set of $t$ distinct coordinates
	there is a \emph{linear} interpolation map $\Sigma^t\to\code$.
\end{definition}

\begin{definition}[Axis-parallel line --- \texttt{def:axis-line}]\label{def:axis-line}
	An axis-parallel line in $[n]^m$ is obtained by fixing all but one
	coordinate.
\end{definition}

\begin{definition}[Tensor code --- \texttt{def:tensor-code}]\label{def:tensor-code}
	\uses{def:axis-line, def:code}
	$\code^{\otimes m}=\{c\colon[n]^m\to\Sigma : c|_\ell\in\code\text{ for every
			axis-parallel line }\ell\}$.
\end{definition}

\begin{proposition}[Tensor-code distance --- \texttt{prop:distance}]\label{prop:distance}
	\uses{def:tensor-code, prop:distance0}
	Two distinct $c,c'\in\code^{\otimes m}$ agree on at most
	$\gamma_m=1-d^m/n^m\le mt/n$ fraction of coordinates.
\end{proposition}

\begin{proposition}[Tuple-to-code correspondence --- \texttt{prop:tuple-to-code-correspondence}]
	\label{prop:tuple-to-code-correspondence}
	\uses{def:tensor-code, prop:distance0}
	Given $t$ distinct coordinates $x_1,\dots,x_t\in[n]$, a tuple
	$(g_1,\dots,g_t)\in(\code^{\otimes m})^t$ comes from at most one
	$h\in\code^{\otimes(m+1)}$ with $h|_{x_j}=g_j$.
\end{proposition}

\begin{proposition}[Interpolation of tuples --- \texttt{prop:interpolate-tuple}]
	\label{prop:interpolate-tuple}
	\uses{def:interpolable, prop:tuple-to-code-correspondence}
	If $\code$ is interpolable then every tuple
	$(g_1,\dots,g_t)\in(\code^{\otimes m})^t$ extends to a codeword in
	$\code^{\otimes(m+1)}$.
\end{proposition}

%% ====================================================================
\section{Section~3: The tensor code test}
\label{sec:test}

\begin{definition}[Tracial strategy --- \texttt{def:tracial-strat}]\label{def:tracial-strat}
	\uses{def:tracial-state, def:submeasurement}
	A tracial strategy $\mathcal{S}=(\tau,A,B,P)$ for the tensor code test
	consists of a tracial state $\tau$ on a von Neumann algebra $\algebra$,
	together with projective measurements $A^u$ (points), $B^\ell$ (lines),
	$P^{u,v}$ (pairs), all taking values in~$\algebra$.
\end{definition}

\begin{definition}[$(\eps,\delta)$-good strategy --- \texttt{def:tracial-good}]
	\label{def:tracial-good}
	\uses{def:tracial-strat, def:consistency, def:processed}
	$\mathcal{S}$ is $(\eps,\delta)$-good if:
	\begin{enumerate}
		\item (Points--lines consistency)
		      $B^\ell_{[g\mapsto g(u)\mid a]}\simeq_\eps A^u_a$
		      on average over a uniformly random axis-parallel line $\ell$ and
		      a uniformly random point $u\in\ell$.
		\item (Points--pairs consistency)
		      $P^{u,v}_{[(a,b)\mapsto a\mid a]}\simeq_\delta A^u_a$
		      and
		      $P^{u,v}_{[(a,b)\mapsto b\mid b]}\simeq_\delta A^v_b$
		      on average over a random subcube $H$ and uniformly random $u,v\in H$.
	\end{enumerate}
\end{definition}

\begin{remark}[Test--goodness bookkeeping]\label{rmk:test-goodness}
	\uses{def:tracial-good}
	A strategy that is $(\eps,\delta)$-good passes the tensor code test with
	probability $\ge 1-\frac12(\eps+\delta)$.
	Conversely, passing with probability $1-\eta$ gives a
	$(2\eta,4\eta)$-good strategy.
\end{remark}

%% ====================================================================
\section{Appendix~A: Operators on expander graphs}
\label{sec:appendix}

\begin{proposition}[Laplacian as an edge average]\label{prop:laplacian-edge-form}
	Let $G$ be a finite graph with a probability distribution on edges
	whose vertex marginals are uniform. Define
	$K=\E_{(u,v)\sim G}|u\rangle\langle v|$ and $L=\frac{1}{n}\id-K$. Then
	$L=\frac12\E_{(u,v)\sim G}(|u\rangle-|v\rangle)(\langle u|-\langle v|)$.
\end{proposition}

\begin{lemma}[Local-to-global expander inequality --- \texttt{lem:local-to-global-expander}]
	\label{lem:local-to-global-expander}
	\uses{prop:laplacian-edge-form}
	Let $G$ be as above, let $\lambda_2$ be the second eigenvalue of the
	normalized Laplacian, let $\{A^u\}_{u\in[n]}$ be bounded operators on a
	Hilbert space $\hilb$, and let $\rho$ be a positive linear functional on
	$\mathcal{B}(\hilb)$. Then
	\[
		\E_{u,v\sim[n]}\rho((A^u-A^v)^*(A^u-A^v))
		\le
		\frac{1}{n\lambda_2}
		\E_{(u,v)\sim G}\rho((A^u-A^v)^*(A^u-A^v)).
	\]
\end{lemma}

%% ====================================================================
\section{Section~5: Self-improvement}
\label{sec:self-improvement}

The self-improvement lemma (\Cref{lem:self-improvement}) is the first of the
two main engines. It relies on three ingredients: orthogonalization, duality
(complementary slackness), and a local-to-global variance bound.

\subsection{Ingredient 1: Orthogonalization}

\begin{lemma}[Orthogonalization --- \texttt{lem:projectivization}]
	\label{lem:projectivization}
	\uses{def:submeasurement, def:tracial-state}
	\notready
	Let $\algebra$ be a von Neumann algebra with a normal state $\tau$ (not
	necessarily tracial). Let $A=\{A_a\}$ be a submeasurement in~$\algebra$
	with $\sum_a\tau(A_a(\id-A_a))\le\zeta$. Then there exists a projective
	submeasurement $\{P_a\}\subset\algebra$ with
	$A_a\approx_{\sqrt{18\zeta}}P_a$.
\end{lemma}

\begin{remark}[External dependency: de Oliveira--Scholz]\label{rmk:projectivization-dep}
	\Cref{lem:projectivization} is proved in the paper by reducing to
	\cite[Theorem~1.2]{de2021orthogonalization} (de Oliveira--Scholz, 2021).
	The proof uses the polar decomposition in a von Neumann algebra.
	This external result must be formalized or axiomatized.
\end{remark}

\subsection{Ingredient 2: Complementary slackness (duality)}

\begin{lemma}[Positive-map duality --- \texttt{lem:duality}]
	\label{lem:duality}
	\uses{def:tracial-state}
	\notready
	Let $\algebra$ be a von Neumann algebra and let
	$\varphi_1,\dots,\varphi_k\colon\algebra\to\C$ be normal positive linear
	maps. Then there is a unique normal linear map $\psi\colon\algebra\to\C$ of
	minimal norm such that $\psi\ge\varphi_i$ for all~$i$. Moreover, there
	exists a measurement $\{T_i\}_{i\in[k]}$ such that for all $X\in\algebra$,
	\[
		\psi(X)=\sum_i\varphi_i(T_i X).
	\]
\end{lemma}

\begin{remark}[External dependency: de Oliveira--Scholz]\label{rmk:duality-dep}
	\Cref{lem:duality} is \cite[Proposition~7.1]{de2021orthogonalization}.
	The proof uses the Hahn--Banach theorem for operator systems.
	This external result must be formalized or axiomatized.
	In finite dimensions, this corresponds to complementary slackness of a
	particular SDP (cf.\ \cite[Lemma~10]{natarajan2018two}), but the paper
	applies it in the infinite-dimensional (von Neumann algebra) setting.
\end{remark}

\subsection{Ingredient 3: Variance bounds}

\begin{lemma}[Local variance along a random line --- \texttt{lem:local-variance}]
	\label{lem:local-variance}
	\uses{def:tracial-good, lem:consistency-consequences, prop:distance}
	\notready
	Let $(\tau,A,B,P)$ be an $(\eps,\delta)$-good strategy and let $\{T_g\}$ be a
	measurement with outcomes in $\code^{\otimes m}$. Then on average over
	axis-parallel lines $\ell$ and points $x,y\in\ell$,
	\[
		A^x_{g(x)}T_g^{1/2}\approx_{\zeta_{\mathrm{local}}}
		A^y_{g(y)}T_g^{1/2},
		\quad
		\zeta_{\mathrm{local}}=2\sqrt{2\eps}+2\sqrt{\gamma},
		\quad
		\gamma=1-d/n.
	\]
\end{lemma}

\begin{lemma}[Global variance --- \texttt{lem:variance}]
	\label{lem:variance}
	\uses{lem:local-variance, lem:local-to-global-expander}
	\notready
	Under the same hypotheses, on average over uniformly random
	$x,y\in[n]^m$,
	\[
		A^x_{g(x)}T_g^{1/2}\approx_{\zeta_{\mathrm{var}}}
		A^y_{g(y)}T_g^{1/2},
		\quad
		\zeta_{\mathrm{var}}=\sqrt{m} \zeta_{\mathrm{local}}.
	\]
	The proof applies \Cref{lem:local-to-global-expander} to the graph on $[n]^m$
	connecting points that differ in at most one coordinate, with
	$\lambda_2=1/(mn^m)$.
\end{lemma}

\subsection{Self-improvement proof}

The self-improvement proof in Section~5 of the paper has four sub-steps.
All labels below are internal to the paper.

\begin{lemma}[Technical helper --- \texttt{lem:self-improvement-helper}]
	\label{lem:self-improvement-helper}
	\uses{lem:variance, prop:holder}
	\notready
	Defines $H^x_g=A^x_{g(x)}\cdot T_g\cdot A^x_{g(x)}$ and
	$H_g=\E_x H^x_g$. Shows
	$\E_x\sum_g\tau(|H_g-A^x_{g(x)}T_gA^x_{g(x)}|)\le 2\zeta_{\mathrm{var}}$.
\end{lemma}

\begin{lemma}[Consistency of $H$ with $A$ --- \texttt{lem:si-cons}]
	\label{lem:si-cons}
	\uses{lem:self-improvement-helper}
	\notready
	On average over $x\sim[n]^m$,
	$H_{[g\mapsto g(x)\mid a]}\simeq_{2\zeta_{\mathrm{var}}} A^x_a$.
\end{lemma}

\begin{lemma}[Projectivity of $H$ --- \texttt{lem:si-proj}]
	\label{lem:si-proj}
	\uses{lem:self-improvement-helper, lem:si-cons, lem:variance, prop:distance}
	\notready
	$\sum_g\tau(H_g-H_g^2)\le 5\zeta_{\mathrm{var}}+\gamma_m$.
\end{lemma}

\begin{lemma}[Self-improvement --- \texttt{lem:self-improvement}]
	\label{lem:self-improvement}
	\uses{lem:projectivization, lem:duality, lem:variance,
		lem:self-improvement-helper, lem:si-cons, lem:si-proj,
		lem:closeness-to-close-ips, lem:transfer-cons, lem:cons-sub-meas}
	\notready
	There exists $\zeta=\operatorname{poly}(m,t)\cdot\operatorname{poly}(\eps,n^{-1})$
	such that:
	given an $(\eps,\delta)$-good strategy and a measurement $\{G_g\}$ that is
	$\nu$-consistent with~$A$, there exists a projective submeasurement $\{H_h\}$ with
	\begin{enumerate}
		\item (Completeness) $\tau(H)\ge 1-\nu-\zeta$,
		\item (Consistency) $H_{[h\mapsto h(x)\mid a]}\simeq_\zeta A^x_a$ on average,
		\item (Boundedness) There exists a positive linear map $\psi$ with
		      $\psi(\id-H)\le\zeta$ and
		      $\psi(X)\ge\E_x\tau(X\cdot A^x_{h(x)})$ for all positive $X$.
	\end{enumerate}
	The proof:
	\begin{enumerate}[label=Step~\arabic*:]
		\item Apply \Cref{lem:duality} to obtain the measurement $\{T_g\}$ and the
		      dominating functional $\psi$.
		\item Define $H_g=\E_x A^x_{g(x)}T_g A^x_{g(x)}$.
		\item Prove completeness, consistency, and projectivity of this (non-projective) $H$
		      (\Cref{lem:si-cons,lem:si-proj}).
		\item Apply \Cref{lem:projectivization} to make $H$ projective.
	\end{enumerate}
\end{lemma}

%% ====================================================================
\section{Section~6: Pasting}
\label{sec:pasting}

Pasting takes ``slice'' submeasurements $\{G^x_g\}_{g\in\code^{\otimes m}}$
for $x\in[n]$ and produces a single measurement with outcomes in
$\code^{\otimes(m+1)}$. The paper gives two methods; Method~2 (with the
Chernoff-based completeness bound) is the one used in the final proof.

\subsection{Commutativity chain}

\begin{lemma}[Commutativity of $A$'s --- \texttt{lem:a-comm}]
	\label{lem:a-comm}
	\uses{def:tracial-good, lem:consistency-consequences, lem:add-a-proj}
	\notready
	On average over uniform $(u,x),(v,y)\sim[n]^m\times[n]$,
	\[
		A^{u,x}_a A^{v,y}_b\approx_{\sqrt{32(m+1)\delta}} A^{v,y}_b A^{u,x}_a.
	\]
\end{lemma}

\begin{lemma}[Evaluated $G$'s commute --- \texttt{lem:g-comm-after-eval}]
	\label{lem:g-comm-after-eval}
	\uses{lem:a-comm, lem:cons-sub-meas, lem:closeness-to-close-ips}
	\notready
	With $\nu_1=8(\sqrt\zeta+\sqrt{(m+1)\delta})$,
	the evaluated operators $G^{u,x}_a=G^x_{[g\mapsto g(u)\mid a]}$ and
	$G^{v,y}_b$ approximately commute on average:
	$G^{u,x}_a G^{v,y}_b\approx_{\nu_1} G^{v,y}_b G^{u,x}_a$.
\end{lemma}

\begin{lemma}[Subspace measurements commute --- \texttt{lem:g-comm}]
	\label{lem:g-comm}
	\uses{lem:g-comm-after-eval, prop:distance}
	\notready
	With $\nu_2=4(\gamma_m+\nu_1)$, on average over $x,y\sim[n]$:
	$G^x_g\,G^y_h\approx_{\nu_2}G^y_h\,G^x_g$.
\end{lemma}

\begin{lemma}[Commutativity switcheroo --- \texttt{lem:commutativity-switcheroo}]
	\label{lem:commutativity-switcheroo}
	\uses{lem:g-comm}
	\notready
	Commutativity of $G^x_g$ with $G^y_h$ implies commutativity with $G^y=\sum_h G^y_h$.
\end{lemma}

\begin{corollary}[Commutativity of $G$'s --- \texttt{cor:g-comm}]
	\label{cor:g-comm}
	\uses{lem:commutativity-switcheroo, lem:g-comm}
	\notready
	Full approximate commutativity: $G^x_g G^y\approx G^y G^x_g$.
\end{corollary}

\begin{corollary}[Commutativity of $\widehat{G}$ --- \texttt{cor:G-hat-facts}]
	\label{cor:G-hat-facts}
	\uses{cor:g-comm, lem:g-comm}
	\notready
	With $\nu'_2=27\nu_2^{1/4}$, the completed measurements
	$\widehat{G}^x_g$ (including the $\bot$ outcome) satisfy
	$\widehat{G}^x_g\,\widehat{G}^y_h\approx_{\nu'_2}\widehat{G}^y_h\,\widehat{G}^x_g$.
\end{corollary}

\subsection{Pasting construction (Method~2)}

Method~2 defines the sandwiched measurement
$\widehat{H}^{x_1,\dots,x_k}_{g_1,\dots,g_k}
	=\widehat{G}^{x_1}_{g_1}\cdots\widehat{G}^{x_k}_{g_k}\cdots\widehat{G}^{x_1}_{g_1}$,
interpolates tuples of type $|\tau|\ge t$ into a global codeword, and averages
over distinct $k$-tuples $(x_1,\dots,x_k)\in\mathrm{distinct}([n],k)$.

\subsection{Consistency of $H$}

\begin{lemma}[Consistency of $\widehat{H}$ with $B$ at one point ---
		\texttt{lem:ld-sandwich-line-one-point}]
	\label{lem:ld-sandwich-line-one-point}
	\uses{cor:G-hat-facts, lem:switcheroo}
	\notready
	Sandwiched products are approximately consistent with the line measurement
	at a single evaluation point.
\end{lemma}

\begin{lemma}[Consistency of $H$ with $B$ --- \texttt{lem:h-b-consistency}]
	\label{lem:h-b-consistency}
	\uses{lem:ld-sandwich-line-one-point, prop:interpolate-tuple,
		prop:tuple-to-code-correspondence, lem:data-processing, lem:transfer-cons}
	\notready
	With error parameter $\nu'_3$, the pasted measurement $H$ is approximately
	consistent with the line measurement~$B$.
\end{lemma}

\begin{corollary}[Consistency of $H$ with $A$ --- \texttt{cor:ha-cons}]
	\label{cor:ha-cons}
	\uses{lem:h-b-consistency, lem:data-processing, lem:transfer-cons}
	\notready
	On average over $(u,x)\in[n]^m\times[n]$,
	$H_{[h\mapsto h(u,x)\mid a]}\simeq_{\nu'_4} A^{u,x}_a$.
\end{corollary}

\subsection{Completeness of $H$}

\begin{lemma}[Over all outcomes --- \texttt{lem:over-all-outcomes}]
	\label{lem:over-all-outcomes}
	\uses{prop:tuple-to-code-correspondence, prop:distance,
		lem:h-b-consistency, cor:G-hat-facts}
	\notready
	With error $\nu'_5$:
	\[
		\tau(H)\approx_{\nu'_5}
		\E_{x_1,\dots,x_k}\sum_{\tau:|\tau|\ge t}
		\sum_{(g_1,\dots,g_k)\in\mathrm{Outcomes}_\tau}
		\tau\!(\widehat{H}^{x_1,\dots,x_k}_{g_1,\dots,g_k}).
	\]
\end{lemma}

\begin{lemma}[From $H$ to $G$ polynomial --- \texttt{lem:from-H-to-G}]
	\label{lem:from-H-to-G}
	\uses{cor:G-hat-facts, lem:switcheroo, lem:closeness-to-close-ips}
	\notready
	With error $\nu'_6=2k^2\nu'_2$:
	\[
		\E_{x_1,\dots,x_k}\sum_{\tau:|\tau|\ge t}
		\sum_{(g_1,\dots,g_k)\in\mathrm{Outcomes}_\tau}
		\tau\!(\widehat{H}^{x_1,\dots,x_k}_{g_1,\dots,g_k})
		\approx_{\nu'_6}
		\sum_{i=t}^k\binom{k}{i}\tau\!(G^i(\id-G)^{k-i}).
	\]
\end{lemma}

\begin{lemma}[Matrix Chernoff for Bernoulli polynomial --- \texttt{lem:chernoff-bernoulli-matrix}]
	\label{lem:chernoff-bernoulli-matrix}
	\uses{def:tracial-state}
	\notready
	Let $F(x)=\sum_{r=t}^k\binom{k}{r}x^r(1-x)^{k-r}$. If
	$0\le X\le\id$ is self-adjoint in~$\algebra$ with $\tau(X)\ge 1-\kappa$,
	then for $k\ge 2t/\theta$ and $0<\theta<1$:
	\[
		\tau(F(X))\ge 1-\frac\kappa{1-\theta}-e^{-\theta^2 k/2}.
	\]
	The proof uses the spectral theorem (functional calculus for bounded
	self-adjoint operators in a von Neumann algebra) and the scalar Chernoff bound.
\end{lemma}

\begin{corollary}[Completeness of $H$ --- \texttt{cor:ld-pasting-N-completeness}]
	\label{cor:ld-pasting-N-completeness}
	\uses{lem:over-all-outcomes, lem:from-H-to-G, lem:chernoff-bernoulli-matrix}
	\notready
	For $k\ge 12mt$ and $\theta=1/(6m)$:
	\[
		\tau(H)\ge 1-\kappa\!\left(1+\frac1{3m}\right)-\nu-e^{-k/(72m^2)},
	\]
	where $\nu=\nu'_5+\nu'_6=\operatorname{poly}(m,t,k)\cdot\operatorname{poly}(\eps,\delta,\zeta,n^{-1})$.
\end{corollary}

\subsection{The pasting lemma (combined)}

\begin{lemma}[Pasting --- \texttt{lem:pasting}]
	\label{lem:pasting}
	\uses{cor:ha-cons, cor:ld-pasting-N-completeness}
	\notready
	There exists a function
	$\mu(\kappa,m,t,r,\eps,\delta,\zeta,n^{-1})
		=\kappa(1+\frac{1}{3m})+\operatorname{poly}(m,t,r)\cdot\operatorname{poly}(\eps,\delta,\zeta,n^{-1})+e^{-r/(72m^2)}$
	such that: given an $(\eps,\delta)$-good strategy for $\code^{\otimes(m+1)}$
	and slice submeasurements $\{G^x_g\}$ satisfying completeness, consistency,
	and boundedness conditions, there exists a pasted measurement $\{H_h\}$
	with outcomes in $\code^{\otimes(m+1)}$ satisfying
	$H_{[h\mapsto h(u,x)\mid a]}\simeq_\mu A^{u,x}_a$ on average.
\end{lemma}

%% ====================================================================
\section{Section~4: Main induction and main theorems}
\label{sec:main-loop}

\begin{lemma}[Induction lemma --- \texttt{lem:induction}]
	\label{lem:induction}
	\uses{lem:self-improvement, lem:pasting}
	\notready
	There exists
	$\sigma(m,t,r,\eps,\delta,n^{-1})=\operatorname{poly}(m,t,r)\cdot(\operatorname{poly}(\eps,\delta,n^{-1})+e^{-\Omega(r/m^2)})$
	such that for any $(\eps,\delta)$-good strategy and $r\ge 12mt$,
	there exists a measurement $\{G_c\}\subset\algebra$ with outcomes in
	$\code^{\otimes m}$ satisfying
	\[
		\E_{u\sim[n]^m}\sum_c\tau(G_c\,A^u_{c(u)})\ge 1-\sigma.
	\]
	Proof is by induction on $m$:
	\begin{itemize}
		\item \textbf{Base case} ($m=1$): the line measurement $B^\ell$ directly
		      plays the role of the global measurement.
		\item \textbf{Inductive step} ($m\to m+1$): apply the induction hypothesis
		      to each slice, apply \Cref{lem:self-improvement} to each slice
		      measurement, then apply \Cref{lem:pasting} to produce the $(m+1)$-dimensional
		      measurement.
	\end{itemize}
\end{lemma}

\begin{theorem}[Quantum soundness --- \texttt{thm:main} (Theorem~4.1)]
	\label{thm:main}
	\uses{lem:induction, lem:self-improvement}
	\notready
	There exists
	$\eta(m,t,r,\eps,n^{-1})=\operatorname{poly}(m,t,r)\cdot\operatorname{poly}(\eps,n^{-1},e^{-\Omega(r/m^2)})$
	such that: if $\mathcal{S}=(\tau,A,B,P)$ is a synchronous strategy passing
	the tensor code test for $\code^{\otimes m}$ with probability $1-\eps$ and
	$r\ge 12mt$, then there exists a projective measurement
	$\{G_c\}_{c\in\code^{\otimes m}}\subset\algebra$ with
	\[
		\E_{u\sim[n]^m}\sum_c\tau(G_c\,A^u_{c(u)})\ge 1-\eta.
	\]

	\textbf{This is the primary formalization target.}
\end{theorem}

\subsection{Extension to general two-prover strategies (Section~4.1)}

\begin{definition}[Two-prover strategy]\label{def:two-prover-strat}
	A two-prover strategy
	$\mathcal{S}=(\ket{\psi},A,B,P,\tilde A,\tilde B,\tilde P)$ with
	$\ket{\psi}\in\hilb_A\otimes\hilb_B$
	($\hilb_A,\hilb_B$ finite-dimensional)
	and separate projective measurements for each prover.
\end{definition}

\begin{definition}[$(\eps,\delta,\xi)$-good two-prover strategy --- \texttt{def:good-2}]
	\label{def:good-2}
	\uses{def:two-prover-strat}
	\notready
	The strategy passes:
	(1)~the axis-parallel line test with probability $\ge 1-\eps$,
	(2)~the subcube commutation test with probability $\ge 1-\delta$,
	(3)~the synchronicity test with probability $\ge 1-\xi$.
\end{definition}

\begin{theorem}[Two-prover extension --- \texttt{thm:main-bipartite} (Theorem~4.8)]
	\label{thm:main-bipartite}
	\uses{thm:main, def:good-2}
	\notready
	There exists
	$\eta'(m,t,k,\eps,n^{-1})=\operatorname{poly}(m,t,k)\cdot\operatorname{poly}(\eps,n^{-1},e^{-k/m^2})$
	such that: if $\mathcal{S}$ is a two-prover strategy passing the two-prover
	tensor code test with probability $1-\eps$ and $k\ge 12mt$, then there
	exist projective measurements $\{G_c\}$ and $\{\tilde G_c\}$ on
	$\hilb_A$ and $\hilb_B$ respectively, satisfying
	\begin{align*}
		A^u_a\otimes\id                        & \simeq_{\eta'}\id\otimes\tilde G_{[c\mapsto c(u)\mid a]}, \\
		G_{[c\mapsto c|_\ell\mid f]}\otimes\id & \simeq_{\eta'}\id\otimes\tilde B^\ell_f,                  \\
		G_c\otimes\id                          & \simeq_{\eta'}\id\otimes\tilde G_c.
	\end{align*}

	\textbf{This is the secondary formalization target.}
\end{theorem}

\begin{remark}[External dependency: vidick2021almost]\label{rmk:bipartite-dep}
	The proof of \Cref{thm:main-bipartite} relies on
	\cite{vidick2021almost} (Vidick, ``Almost synchronous quantum correlations'',
	2021) to reduce from general two-prover strategies to the synchronous
	setting of \Cref{thm:main}. This external result must be formalized or
	axiomatized.
\end{remark}

%% ====================================================================
\section{Dependency summary}
\label{sec:deps}

The proof DAG in linearized order (paper labels):

\small
\begin{longtable}{L{0.04\textwidth}L{0.32\textwidth}L{0.28\textwidth}L{0.26\textwidth}}
	\toprule
	\# & Label                                      & Paper location & Direct dependencies                     \\
	\midrule
	\endfirsthead
	\toprule
	\# & Label                                      & Paper location & Direct dependencies                     \\
	\midrule
	\endhead
	1  & \texttt{def:tracial-state}                 & Sec.~2.1       & ---                                     \\
	2  & \texttt{def:tau-norm}                      & Sec.~2.1       & 1                                       \\
	3  & \texttt{prop:holder}                       & Sec.~2.1       & 1, 2                                    \\
	4  & \texttt{cor:cauchy-schwarz}                & Sec.~2.1       & 1                                       \\
	5  & \texttt{def:submeasurement}                & Sec.~2.2       & 1                                       \\
	6  & \texttt{def:processed}                     & Sec.~2.2       & 5                                       \\
	7  & \texttt{def:consistency}                   & Sec.~2.2       & 2, 5                                    \\
	8  & \texttt{lem:data-processing}               & Sec.~2.2       & 6, 7                                    \\
	9  & \texttt{lem:consistency-consequences}      & Sec.~2.2       & 7                                       \\
	10 & \texttt{lem:closeness-to-consistency}      & Sec.~2.2       & 7                                       \\
	11 & \texttt{lem:closeness-to-close-ips}        & Sec.~2.2       & 7, 4                                    \\
	12 & \texttt{lem:add-a-proj}                    & Sec.~2.2       & 7                                       \\
	13 & \texttt{lem:transfer-cons}                 & Sec.~2.2       & 7, 11                                   \\
	14 & \texttt{lem:cons-sub-meas}                 & Sec.~2.2       & 7                                       \\
	15 & \texttt{lem:switcheroo}                    & Sec.~2.2       & 12                                      \\
	16 & \texttt{def:code}                          & Sec.~2.3       & ---                                     \\
	17 & \texttt{prop:distance0}                    & Sec.~2.3       & 16                                      \\
	18 & \texttt{def:interpolable}                  & Sec.~2.3       & 16, 17                                  \\
	19 & \texttt{def:axis-line}                     & Sec.~2.3       & ---                                     \\
	20 & \texttt{def:tensor-code}                   & Sec.~2.3       & 16, 19                                  \\
	21 & \texttt{prop:distance}                     & Sec.~2.3       & 17, 20                                  \\
	22 & \texttt{prop:tuple-to-code-correspondence} & Sec.~2.3       & 17, 20                                  \\
	23 & \texttt{prop:interpolate-tuple}            & Sec.~2.3       & 18, 22                                  \\
	24 & \texttt{def:tracial-strat}                 & Sec.~3         & 1, 5                                    \\
	25 & \texttt{def:tracial-good}                  & Sec.~3         & 6, 7, 24                                \\
	26 & \texttt{prop:laplacian-edge-form}          & App.~A         & ---                                     \\
	27 & \texttt{lem:local-to-global-expander}      & App.~A         & 26                                      \\
	28 & \texttt{lem:projectivization}              & Sec.~5         & 1, 5; \textbf{ext:}~de Oliveira--Scholz \\
	29 & \texttt{lem:duality}                       & Sec.~5         & 1; \textbf{ext:}~de Oliveira--Scholz    \\
	30 & \texttt{lem:local-variance}                & Sec.~5         & 9, 21, 25                               \\
	31 & \texttt{lem:variance}                      & Sec.~5         & 27, 30                                  \\
	32 & \texttt{lem:self-improvement-helper}       & Sec.~5         & 3, 31                                   \\
	33 & \texttt{lem:si-cons}                       & Sec.~5         & 32                                      \\
	34 & \texttt{lem:si-proj}                       & Sec.~5         & 21, 31, 32, 33                          \\
	35 & \texttt{lem:self-improvement}              & Sec.~5         & 28, 29, 31, 32, 33, 34, 11, 13, 14      \\
	36 & \texttt{lem:a-comm}                        & Sec.~6.1       & 9, 12, 25                               \\
	37 & \texttt{lem:g-comm-after-eval}             & Sec.~6.2       & 14, 36                                  \\
	38 & \texttt{lem:g-comm}                        & Sec.~6.2       & 21, 37                                  \\
	39 & \texttt{lem:commutativity-switcheroo}      & Sec.~6.4       & 38                                      \\
	40 & \texttt{cor:g-comm}                        & Sec.~6.4       & 38, 39                                  \\
	41 & \texttt{cor:G-hat-facts}                   & Sec.~6.4       & 38, 40                                  \\
	42 & \texttt{lem:ld-sandwich-line-one-point}    & Sec.~6.5       & 15, 41                                  \\
	43 & \texttt{lem:h-b-consistency}               & Sec.~6.5       & 8, 13, 22, 23, 42                       \\
	44 & \texttt{cor:ha-cons}                       & Sec.~6.5       & 8, 13, 43                               \\
	45 & \texttt{lem:over-all-outcomes}             & Sec.~6.6       & 21, 22, 41, 43                          \\
	46 & \texttt{lem:from-H-to-G}                   & Sec.~6.6       & 11, 15, 41                              \\
	47 & \texttt{lem:chernoff-bernoulli-matrix}     & Sec.~6.6       & 1 (functional calculus)                 \\
	48 & \texttt{cor:ld-pasting-N-completeness}     & Sec.~6.6       & 45, 46, 47                              \\
	49 & \texttt{lem:pasting}                       & Sec.~6.7       & 44, 48                                  \\
	50 & \texttt{lem:induction}                     & Sec.~4         & 35, 49                                  \\
	51 & \texttt{thm:main}                          & Sec.~4         & 35, 50                                  \\
	52 & \texttt{def:good-2}                        & Sec.~4.1       & 24                                      \\
	53 & \texttt{thm:main-bipartite}                & Sec.~4.1       & 51, 52; \textbf{ext:}~vidick2021almost  \\
	\bottomrule
\end{longtable}
\normalsize

%% ====================================================================
\section{External dependencies}
\label{sec:external}

The paper directly invokes the following external results. They must either be
formalized or introduced as axioms with careful type signatures.

\subsection{de Oliveira--Scholz (2021)}

Reference: \cite{de2021orthogonalization}
(``Orthogonalization of positive operator valued measures'').

\begin{enumerate}
	\item \textbf{Theorem~1.2} (orthogonalization): used for
	      \texttt{lem:projectivization} (\Cref{lem:projectivization}).
	      Given a measurement on a von Neumann algebra with a normal state that is
	      almost projective, produces a nearby projective submeasurement.
	      The proof uses polar decomposition in a von Neumann algebra.

	\item \textbf{Proposition~7.1} (positive-map duality): used for
	      \texttt{lem:duality} (\Cref{lem:duality}).
	      Given finitely many normal positive linear maps on a von Neumann algebra,
	      produces a minimal dominating normal positive linear map with a
	      complementary-slackness decomposition.
	      The proof uses the Hahn--Banach theorem for operator systems.
\end{enumerate}

\subsection{Vidick (2021)}

Reference: \cite{vidick2021almost}
(``Almost synchronous quantum correlations'').

Used in the proof of \texttt{thm:main-bipartite}
(\Cref{thm:main-bipartite}) to reduce general two-prover strategies to
the synchronous setting. The specific results used are the ``rounding''
lemmas that produce near-synchronous behavior from a near-perfect
synchronicity score.

\subsection{Other background}

\begin{itemize}
	\item \textbf{Spectral theorem / functional calculus} for bounded
	      self-adjoint operators on a Hilbert space: used in the proof of
	      \texttt{lem:chernoff-bernoulli-matrix} (\Cref{lem:chernoff-bernoulli-matrix}).
	\item \textbf{Hahn--Banach theorem}: used indirectly via
	      \cite{de2021orthogonalization}; see above.
	\item \textbf{Polar decomposition in a von Neumann algebra}: used
	      indirectly via \cite{de2021orthogonalization}; see above.
	\item \textbf{Additive Chernoff bound}: used inside the proof of
	      \texttt{lem:chernoff-bernoulli-matrix} after reducing to scalar eigenvalues.
\end{itemize}

%% ====================================================================
\section{Non-goals (substitutes explicitly excluded)}
\label{sec:nongoals}

The following approaches are \textbf{not} faithful to the paper's proof and are
explicitly excluded from this blueprint. They may be useful for pedagogy or for
a finite-dimensional pilot, but they are not substitutes for the strict plan.

\begin{enumerate}
	\item \textbf{Finite-dimensional-only pilot replacements.}
	      The paper's ambient setting is a von Neumann algebra with a normal tracial
	      state, not $\Mat_D(\C)$ with normalized trace.
	      All definitions, lemma statements, and proofs in this blueprint use the
	      paper's full generality. A finite-dimensional specialization may be pursued
	      separately as a warm-up, but it does not discharge the nodes in this blueprint.

	\item \textbf{SDP reformulations replacing \texttt{lem:duality}.}
	      The paper's duality lemma (\Cref{lem:duality}) is a result about normal
	      positive linear maps on a von Neumann algebra, proved via Hahn--Banach.
	      Replacing it with finite-dimensional SDP strong duality (as in
	      \cite[Lemma~10]{natarajan2018two}) changes the proof's logic and ambient
	      setting.

	\item \textbf{SVD replacements for \texttt{lem:projectivization}.}
	      The paper's orthogonalization lemma (\Cref{lem:projectivization}) uses the
	      polar decomposition in a von Neumann algebra, citing
	      \cite{de2021orthogonalization}.
	      A finite-dimensional SVD-based orthogonalization is a different proof.

	\item \textbf{Importing Paper~2 or Paper~1 machinery.}
	      This blueprint covers arXiv:2111.08131 as a standalone paper.
	      The notation and lemma numbering from the full MIP*=RE proof
	      (arXiv:2001.04383, arXiv:2012.12817) should not be imported unless
	      arXiv:2111.08131 itself directly invokes them.
	      The only external references the paper uses are
	      \cite{de2021orthogonalization} and \cite{vidick2021almost}.
\end{enumerate}

%% ====================================================================
\section{Suggested formalization order}
\label{sec:order}

\subsection{Phase~1: Definitions and Section~2 infrastructure}

Formalize the tracial-state API, measurements, consistency/closeness, all
Section~2 propositions, and the coding-theoretic definitions.

\subsection{Phase~2: Section~3 test interface and Appendix~A}

Formalize the tensor code test, tracial strategies,
$(\eps,\delta)$-good strategies, and the expander lemma.

\subsection{Phase~3: Section~5 self-improvement}

Formalize the variance bounds (\texttt{lem:local-variance},
\texttt{lem:variance}), the self-improvement helper lemmas
(\texttt{lem:self-improvement-helper}, \texttt{lem:si-cons},
\texttt{lem:si-proj}), and the full self-improvement lemma.
This phase requires axiomatizing or formalizing the external dependencies
\texttt{lem:projectivization} and \texttt{lem:duality}.

\subsection{Phase~4: Section~6 pasting}

Formalize the commutativity chain
(\texttt{lem:a-comm} through \texttt{cor:G-hat-facts}),
the pasting construction (Method~2), the consistency analysis
(\texttt{lem:h-b-consistency}, \texttt{cor:ha-cons}),
the completeness analysis
(\texttt{lem:over-all-outcomes}, \texttt{lem:from-H-to-G},
\texttt{lem:chernoff-bernoulli-matrix}, \texttt{cor:ld-pasting-N-completeness}),
and the combined pasting lemma.

\subsection{Phase~5: Section~4 main theorems}

Formalize the induction lemma (\texttt{lem:induction}),
the main theorem (\texttt{thm:main}), the two-prover definitions,
and the bipartite extension (\texttt{thm:main-bipartite}).
The bipartite extension requires axiomatizing or formalizing
\cite{vidick2021almost}.

%% ====================================================================
\begin{thebibliography}{99}
	\bibitem{de2021orthogonalization}
	L.~de Oliveira, V.~Scholz,
	\emph{Orthogonalization of positive operator valued measures},
	2021.

	\bibitem{vidick2021almost}
	T.~Vidick,
	\emph{Almost synchronous quantum correlations},
	2021.

	\bibitem{natarajan2018two}
	A.~Natarajan, T.~Vidick,
	\emph{Two-player entangled games are NP-hard},
	2018.

	\bibitem{blackadar2006operator}
	B.~Blackadar,
	\emph{Operator Algebras: Theory of C*-Algebras and von Neumann Algebras},
	Springer, 2006.

	\bibitem{pisier2003non}
	G.~Pisier,
	\emph{Non-commutative $L^p$ Spaces},
	2003.

	\bibitem{hall2013quantum}
	B.~C.~Hall,
	\emph{Quantum Theory for Mathematicians},
	Springer, 2013.
\end{thebibliography}

\end{document}